% 2de graad, reeks Veeltermen.
% Malevitsj: vierkanten en rechthoeken van verschillende grootte.
% Motief: algebrategels, (x + 1)(x + 2) = x^2 + 3x + 2 als oppervlakte.
\def\omslagbestand{omslag-suprematisme.jpg}
\def\omslagbron{Kazimir Malevitsj, \textit{Suprematisme}, 1915 (publiek domein)}
\def\omslagrand{10.4 + 0.04*\x}
\def\omslagmotief{%
  \randlijn
  \fill[fel] (15,5.2) rectangle (17,7.2);
  \fill[white,fill opacity=.25] (17,5.2) rectangle (18.6,7.2) (15,7.2) rectangle (17,8);
  \fill[white,fill opacity=.55] (17,7.2) rectangle (18.6,8);
  \draw[white,line width=1pt] (15,5.2) rectangle (18.6,8);
  \draw[white,line width=1pt] (17,5.2) -- (17,8) (17.8,5.2) -- (17.8,8) (15,7.2) -- (18.6,7.2);
  \node[text=inkt] at (16,6.2) {$x^2$};
  \node[text=white,font=\small,anchor=north] at (16,5.15) {$x$};
  \node[text=white,font=\small,anchor=north] at (17.4,5.15) {$1$};
  \node[text=white,font=\small,anchor=north] at (18.2,5.15) {$1$};
  \node[text=white,font=\small,anchor=east] at (14.95,6.2) {$x$};
  \node[text=white,font=\small,anchor=east] at (14.95,7.6) {$1$};
  \node[text=white,anchor=south] at (16.8,8.15) {$(x+1)(x+2)=x^2+3x+2$};
}
